\section{Introduction du Chapitre}
Après avoir détaillé la réalisation technique du projet, ce dernier chapitre propose une réflexion sur l'expérience humaine et professionnelle acquise au cours de ce stage chez GIAS Industries. Nous y abordons les défis rencontrés, les solutions apportées ainsi qu'un bilan des apports mutuels entre l'étudiant et l'entreprise d'accueil.

\section{Analyse des Problèmes et Solutions}

\subsection{Défis Techniques}
Le principal défi technique a été l'intégration de services tiers comme Twilio pour l'envoi de SMS OTP, nécessitant une gestion fine des clés API et des environnements de test. De plus, la génération dynamique de documents PDF complexes (conventions de stage) a demandé une maîtrise approfondie de la bibliothèque PDFKit pour assurer une mise en page professionnelle respectant la charte graphique de GIAS.

\subsection{Gestion des Difficultés}
Pour surmonter ces obstacles, nous avons adopté une démarche itérative. Chaque problème a été décomposé en sous-tâches plus simples, et nous avons eu recours à une documentation technique rigoureuse. L'appui des ressources en ligne et les échanges avec les encadrants ont été déterminants pour valider les choix technologiques.

\section{Bilan des Apports du Stage}

\subsection{Apports par rapport à la Formation Théorique}
Ce stage a été une opportunité unique de mettre en pratique les concepts de base de données (SQL), de programmation web (Javascript/Node.js) et de gestion de projet (Agile) étudiés durant notre cursus de Licence en Informatique de Gestion. Il a permis de comprendre que la technique doit toujours être au service d'un besoin métier concret.

\subsection{Apports pour GIAS Industries}
L'application développée apporte une valeur ajoutée immédiate à la DRH en centralisant toutes les données des stagiaires dans un outil unique. Le gain de temps sur la génération des conventions et la fiabilité du suivi des présences permettent une gestion plus sereine et plus transparente du capital humain.

\section{Découverte du Métier de Gestionnaire RH}
L'immersion chez GIAS Industries nous a permis de découvrir les facettes du métier de gestionnaire RH dans une structure industrielle. Nous avons compris l'importance de la rigueur administrative, de la confidentialité des données et de la réactivité nécessaire pour répondre aux besoins des différents départements de l'usine en termes de main-d'œuvre temporaire et de stagiaires.

\section{Conclusion du Chapitre}
L'analyse de ce stage montre que la réussite d'un projet informatique dépend autant de la compréhension du métier que de la maîtrise des technologies. Cette expérience chez GIAS Industries a renforcé notre désir de poursuivre une carrière alliant informatique et gestion d'entreprise.